
\textbf{Requisitos Funcionais}:

\subsubsection{Módulo de Autenticação e Controle de Acesso}
\begin{itemize}
\item RF001 – Permitir login por e-mail e senha.
\item RF002 – Permitir redefinição de senha via e-mail.
\item RF003 – Validar credenciais antes de conceder acesso.
\item RF004 – Permitir cadastro de novos usuários (apenas administrador).
\item RF005 – Permitir gerenciamento de perfis e permissões.
\item RF006 – Controlar acesso a áreas restritas com base no status de autenticação do usuário.
\item RF007 – Controlar sessões ativas por conta, permitindo apenas uma sessão simultânea.
\end{itemize}

\subsubsection{Módulo de Templates de Certificados}
\begin{itemize}
\item RF008 – Permitir upload de templates (PPTX, DOCX, PDF).
\item RF009 – Permitir personalização de templates.
\item RF010 – Permitir selecionar template da galeria interna.
\item RF011 – Identificar campos dinâmicos automaticamente.
\item RF012 – Permitir criação/edição manual de campos dinâmicos (ex: \{\{nome\}\}).
\item RF013 – Salvar configurações de mapeamento de campos.
\item RF014 – Permitir duplicar templates existentes.
\end{itemize}

\subsubsection{Módulo de Entrada de Dados}
\begin{itemize}
\item RF015 – Permitir cadastro individual de participantes.
\item RF016 – Permitir emissão manual para participante único.
\item RF017 – Importar dados via CSV.
\item RF018 – Importar dados via Excel (.xlsx).
\item RF019 – Configurar gatilhos de emissão automática.
\item RF020 – Permitir upload de lista em lote.
\item RF021 – Processar múltiplos certificados em fila.
\item RF022 – Notificar ao término do processamento.
\end{itemize}

\subsubsection{Módulo de Geração de Certificados}
\begin{itemize}
\item RF023 – Gerar certificados em PDF.
\item RF024 – Preencher automaticamente campos dinâmicos.
\item RF025 – Gerar código único por certificado.
\item RF026 – Registrar data, hora e localidade da emissão.
\item RF027 – Permitir reemissão mantendo histórico.
\item RF028 – Permitir download individual.
\item RF029 – Permitir download em lote.
\end{itemize}

\subsubsection{Módulo de Envio e Comunicação}
\begin{itemize}
\item RF030 – Permitir personalização do corpo do e-mail com variáveis.
\item RF031 – Registrar status de envio (pendente, enviado, falha).
\item RF032 – Reenviar automaticamente em caso de falha (até 3 tentativas a cada 10 min).
\end{itemize}

\subsubsection{Módulo de Validação Pública e Antifraude}
\begin{itemize}
\item RF033 – Disponibilizar portal público de validação.
\item RF034 – Validar certificado via código único.
\item RF035 – Gerar QR Code no certificado.
\item RF036 – Permitir personalização de cores do QR Code.
\item RF037 – Exibir informações básicas ao validar (nome, curso, data, validade, autenticidade).
\end{itemize}

\subsubsection{Módulo de Assinaturas Digitais}
\begin{itemize}
\item RF038 – Permitir adicionar assinatura digitalizada ao template.
\item RF039 – Registrar hash de verificação da assinatura.
\end{itemize}

\subsubsection{Módulo de Dashboard e Gestão}
\begin{itemize}
\item RF040 – Exibir lista completa de certificados emitidos.
\item RF041 – Permitir filtros por curso, período, status e instrutor.
\item RF042 – Exibir estatísticas mensais de emissão.
\item RF043 – Exibir logs de ações realizadas.
\end{itemize}

\subsubsection{Módulo de Auditoria e Segurança Operacional}
\begin{itemize}
\item RF044 – Registrar responsável por cada certificado gerado.
\item RF045 – Registrar data, hora e local de ações relevantes.
\item RF046 – Permitir arquivamento de certificados como alternativa à exclusão definitiva.
\end{itemize}

\subsubsection{Módulo de Gestão de Cursos/Eventos}
\begin{itemize}
\item RF047 – Permitir cadastro de cursos/eventos.
\item RF048 – Associar participantes ao curso.
\item RF049 – Definir carga horária.
\item RF050 – Registrar instrutor responsável.
\item RF051 – Configurar datas de início e término.
\item RF052 – Listar certificados vinculados ao curso.
\item RF053 – Validar vínculo do participante ao curso antes de permitir a emissão.
\item RF054 – Registrar local do curso.
\end{itemize}

\subsubsection{Módulo de Participantes}
\begin{itemize}
\item RF055 – Cadastro completo de participantes (nome, e-mail, CPF opcional).
\item RF056 – Permitir edição de dados.
\item RF057 – Importar participantes sem duplicidade.
\item RF058 – Exibir histórico de certificados por participante.
\item RF059 – Permitir busca por nome, e-mail e CPF.
\end{itemize}

\subsubsection{Módulo Multiempresa / Multi-tenant}
\begin{itemize}
\item RF060 – Permitir cadastro de múltiplas organizações.
\item RF061 – Isolar dados entre organizações.
\item RF062 – Permitir que cada empresa tenha seus próprios templates.
\item RF063 – Permitir que administradores gerenciem usuários internos.
\end{itemize}

\subsubsection{Módulo de Configurações do Sistema}
\begin{itemize}
\item RF064 – Personalizar mensagens padrão.
\item RF065 – Ativar/desativar envio automático.
\item RF066 – Configurar regras de reemissão.
\end{itemize}

\subsubsection{Módulo de Armazenamento e Arquivos}
\begin{itemize}
\item RF067 – Permitir exclusão lógica (soft-delete) de arquivos.
\item RF068 – Armazenar templates no sistema.
\end{itemize}

\subsubsection{Módulo de Notificações e Alertas}
\begin{itemize}
\item RF069 – Notificar falhas de emissão.
\item RF070 – Notificar sucesso no processamento de planilhas.
\item RF071 – Alertar quando template estiver incompleto ou inválido.
\end{itemize}

\subsubsection{Módulo de Monitoramento e Logs}
\begin{itemize}
\item RF072 – Registrar logs de autenticação e tentativas de acesso.
\item RF073 – Registrar eventos críticos (falha geração PDF etc.).
\item RF074 – Permitir exportação de logs.
\end{itemize}

\subsubsection{Módulo de Suporte, Operação e Manutenção}
\begin{itemize}
\item RF075 – Permitir abertura de chamados de suporte.
\item RF076 – Exibir status do sistema (online/manutenção).
\item RF077 – Permitir atualização de termos de uso e políticas.
\end{itemize}